\documentclass[twocolumn]{aastex631}
\begin{document}
\title{A residual reionization ionizing-photon-budget shortfall for star-forming galaxies at z\~{}7 (lowxi systematic corner)}
\author{NebulaMind Lab (autonomous pipeline)}
\affiliation{NebulaMind Lab, \url{https://lab.nebulamind.net}}
\begin{abstract}
Reionization ionizing-photon-budget reconciliation at z\~{}7: star-forming galaxies require f\_esc=0.251 (+0.253/-0.128) to close the budget (Madau-Dickinson SFRD, log xi\_ion=25.5+/-0.15, clumping C=2-5, JWST-SFRD tail), versus indirect-proxy-inferred f\_esc=0.062 (+0.108/-0.039) from LzLCS O32/beta calibrations. Median delta(required-inferred)=+0.169 dex-frac (16-84\%: +0.018 to +0.422); 87\% of the systematic MC shows a shortfall. Conclusion: a genuine SHORTFALL remains, and the sign holds under both O32 and beta calibrations. The result is bounded by the xi\_ion x clumping x proxy-calibration systematic, not by statistics. Generated autonomously from public data (jwst) via the NebulaMind Lab runner: a bounded, reproducible, descriptive study.
\end{abstract}
\section{Introduction}
Recent studies have highlighted a potential crisis in our understanding of reionization, with concerns that current models may not produce enough ionizing photons to match observations [Muñoz2024]. This issue is further complicated by the need for accurate estimates of the escape fraction (f\_esc) of ionizing photons from star-forming galaxies, as emphasized by Davies et al. [Davies2021] and Park et al. [Park2022]. To address this challenge, we revisit the reionization-photon-budget using a literature-anchored budget calculation.
\section{Data and method}
Our approach relies on established values from the literature: the cosmic SFRD is based on the Madau \& Dickinson (2014) analytic fitting function, while xi\_ion and O32/beta f\_esc proxy calibrations are adopted from published works [LzLCS: Chisholm+22, Flury+22; Simmonds+24]. We do not utilize survey catalog data or observations from JWST, SDSS, or TNG in this study. Instead, we focus on reconciling systematics across previously published literature values using an ionizing-photon-budget method.
\section{Result}
Our calculation reveals that star-forming galaxies at z\~{}7 require a higher escape fraction (f\_esc=0.251 +0.253/-0.128) to reconcile the reionization photon budget than what is inferred from indirect-proxy methods (f\_esc=0.062 +0.108/-0.039). This discrepancy results in a median shortfall of +0.169 dex-frac, with 87\% of our systematic Monte Carlo simulations showing a deficit. Notably, this result holds under both O32 and beta calibrations.
\begin{figure}\centering\includegraphics[width=\columnwidth]{result.png}\caption{A residual reionization ionizing-photon-budget shortfall for star-forming galaxies at z\~{}7 (lowxi systematic corner)}\end{figure}
\section{Caveats}
It is essential to acknowledge the limitations of our approach. Our calculation relies on automated, single-selection methods without direct calibration against observational data, which may introduce biases or inaccuracies. Furthermore, we depend on previously published values for key parameters like xi\_ion and f\_esc proxy calibrations, which themselves carry uncertainties. These factors restrict the precision and reliability of our findings, emphasizing the need for further research and refined measurements to better understand the reionization photon budget.
\section*{References}
\footnotesize
[Muoz2024] Muñoz et al. (2024). Reionization after JWST: a photon budget crisis?. bibcode 2024MNRAS.535L..37M \par [Davies2021] Davies et al. (2021). The Predicament of Absorption-dominated Reionization: Increased Demands on Ionizing Sources. bibcode 2021ApJ...918L..35D \par [Park2022] Park et al. (2022). Calibrating excursion set reionization models to approximately conserve ionizing photons. bibcode 2022MNRAS.517..192P \par [Duncan2015] Duncan et al. (2015). Powering reionization: assessing the galaxy ionizing photon budget at z < 10. bibcode 2015MNRAS.451.2030D \par [Madau2017] Madau (2017). Cosmic Reionization after Planck and before JWST: An Analytic Approach. bibcode 2017ApJ...851...50M

\end{document}
